\chapter{Cheat Sheet}
\label{app:cheatsheet}

A compact reference to Salam's syntax and the most-used standard-library calls.
Every snippet is valid Salam; refer back to the named chapter for the full story.

\section*{Hello, world (Chapter 1)}

\begin{salam}
func main:
    println "Hello, Salam!"
end
\end{salam}

\noindent Run it: \code{salam exec hello.salam}\quad Build it:
\code{salam build hello.salam --output=hello}

\section*{Variables and types (Chapters 2--3)}

\begin{salam}
name: str = "Sara"          // immutable
mut count: int = 0          // mutable
const MAX = 100             // compile-time constant
x: auto = 3.14              // inferred (f64)
y: int = 250                // int by default
b: u8 = 250 as int as u8    // explicit conversion with 'as'
\end{salam}

\noindent Types: \code{i8 i16 i32 i64}, \code{u8 u16 u32 u64}, \code{f32 f64},
\code{bool}, \code{char}, \code{str}; \code{int}, \code{uint} and \code{float}
are the short spellings of \code{i32}, \code{u32} and \code{f32}. \\
Operators: \code{+ - * / \%}, \code{**} (power, float), \code{== != < > <= >=},
\code{\&\& || !}, \code{+= -= *= /= \%= **=}. \\
Integer \code{/} truncates; no bitwise operators.

\section*{Control flow (Chapter 4)}

\begin{salam}
if x > 0:
    println "positive"
else x == 0:
    println "zero"
else:
    println "negative"
end
\end{salam}

\section*{Loops (Chapter 5)}

\begin{salam}
until cond:           ... end
repeat 3:             ... end          // 3 times
repeat n with i:      ... end          // i = 0..n-1
repeat 1 to 5:        ... end          // 1..5 inclusive
repeat 0 to 20 by 2:      ... end      // step by 2
each x in xs:         ... end          // walk a collection
// break  -- exit loop;  continue  -- next iteration
\end{salam}

\section*{Functions (Chapter 6)}

\begin{salam}
func add(a: int, b: int): int:
    ret a + b
end

func greet(name: str):       // no return type = void
    println "Hi,", name
end
\end{salam}

\noindent Pass by value; no default or reference parameters; overloading by
parameter types; definition order does not matter.

\section*{Arrays and vectors (Chapter 7)}

\begin{salam}
a: int[3] = [1, 2, 3]            // fixed array, indexed 0..2
g: int[2][3] = [[1,2,3],[4,5,6]] // 2-D

mut v: Vector<int> = Vector {}
v.push(10)        // append
v.get(0)[0]       // read element 0
v.set(0, 99)      // write
v.len()           // count
v.pop()           // remove last
v.free()          // release memory
\end{salam}

\section*{Strings (Chapter 8)}

\begin{salam}
import "str"
"a" + "b"                  // concatenation
len("hello")               // 5
str.Upper(s)  str.Lower(s)  str.Trim(s)  str.Reverse(s)
str.Find(s, sub)  str.Contains(s, sub)  str.Replace(s, a, b)
str.Substr(s, start, n)
str.Split(s, sep)  str.Join(vec, sep)
str.FromInt(n as i64)  str.ToInt(s)
\end{salam}

\section*{Structs and methods (Chapter 9)}

\begin{salam}
struct Point:
    x: int = 0
    y: int = 0
    func move_by(dx: int, dy: int):
        this.x = this.x + dx
        this.y = this.y + dy
    end
end

mut p: Point = Point { x = 1, y = 2 }
p.move_by(3, 4)
println p.x, p.y
\end{salam}

\section*{Enums (Chapter 10)}

\begin{salam}
enum Day: Mon, Tue, Wed, Thu, Fri, Sat, Sun end
d: Day = Day.Sat
println d as int           // 5
// "match" with if / else if (no match statement)
\end{salam}

\section*{Interfaces and generics (Chapters 11--12)}

\begin{salam}
interface Shape:
    func area(): f64
end

func describe<T: Shape>(s: T):     // static bound
    println s.area()
end

func draw(s: dyn Shape): ... end   // dynamic dispatch

func Max<T>(a: T, b: T): T:          // generic function
    if a > b: ret a end
    ret b
end

struct Box<T>: value: T end        // generic struct
\end{salam}

\section*{Closures (Chapter 13)}

\begin{salam}
double: auto = (x: int) => x * 2
apply((x: int) => x + 1, 5)            // lambda as argument
func make_adder(n: int): func(int) int:
    ret (x: int) => x + n               // closure capturing n
end
\end{salam}

\section*{Errors and defer (Chapter 14)}

\begin{salam}
func withdraw(n: int): bool: ... end    // bool success flag

import "option"
option.Some(x)  option.IsSome(o)  option.UnwrapOr(o, fallback)
Option { value = 0, present = false } // a "none"

defer v.free()                        // runs at scope exit (LIFO)
\end{salam}

\section*{Modules (Chapter 15)}

\begin{salam}
package main
import "math" strutil "str" "io"      // strutil is an alias for str

pub func public_fn(): ... end         // exported; default is private
\end{salam}

\section*{Standard library (Chapters 16--17)}

\begin{salam}
io.Print(s)  io.Println(s)  io.Input()  io.ReadFile(p)  io.WriteFile(p, t)
fmt.Int(n as i64)  fmt.Float(x)  fmt.PadLeft(s, w)
math.Sqrt  math.Pow  math.Floor  math.MaxI  math.Gcd  math.PI
rand.Seed(s as i64)  rand.IntRange(lo, hi)  rand.Float()  rand.Bool()
time.Now()  time.Format(t)  time.Sleep(ms)
HashMap: insert(k,v)  get(k)  has(k)  remove(k)  size()  free()
Stack: push/pop/peek;  Queue: enqueue/dequeue
json.Valid  json.Get  json.GetInt  json.Has
regex.MatchStr(pat, s)  regex.FindStr(pat, s)  regex.ReplaceStr(pat, s, w)
os.Exists  os.Remove  os.Mkdir  os.ListDir  os.Args  os.Env  os.Exit
\end{salam}

\section*{Memory and FFI (Chapters 18--19)}

\begin{salam}
import "mem"
mem.Allocate(n as u64)  mem.Free(p)  mem.Reallocate  mem.Copy

extern:
    func printf(format: str, ...): int
    func sqrt(x: f64): f64
end
link "sqlite3"                        // link an external library
\end{salam}

\section*{Concurrency (Chapter 20)}

\begin{salam}
import "sync"
mut lock: sync.Mutex = sync.NewMutex()
t: i64 = spawn(work)                   // start a thread
join(t)                               // wait for it
sync.Lock(lock) ... sync.Unlock(lock) // critical section
sync.Destroy(lock)
\end{salam}

\section*{Testing (Chapter 21)}

\begin{salam}
import "testing"
testing.AssertEqInt(got, want)
testing.AssertEqStr(got, want)
testing.AssertTrue(cond)
testing.Summary()                     // prints "tests: N passed, M failed"
\end{salam}

\section*{The \texttt{salam} commands}

\begin{center}
\begin{tabularx}{\linewidth}{@{}lX@{}}
\toprule
\textbf{Command} & \textbf{Does} \\
\midrule
\code{salam exec f.salam}   & run with the interpreter (no C toolchain) \\
\code{salam run f.salam}    & build and run, keep nothing \\
\code{salam build f.salam --output=f} & build a native executable \\
\code{salam fmt f.salam}    & reformat source in place \\
\code{salam new name}       & scaffold a new project \\
\code{salam debug f.salam}  & build with debug info, launch a debugger \\
\code{salam memcheck f.salam} & build with AddressSanitizer and run \\
\code{salam ... --lang=fa}  & treat the source as Persian \\
\bottomrule
\end{tabularx}
\end{center}
